\chapter{Recomendaciones de operacion del torneo}
\label{cap:recomendaciones-torneo}

\begin{procedimiento}{Objetivo del capitulo}{Organizador}
Reducir riesgos antes, durante y despues del evento mediante recomendaciones operativas basadas en los flujos implementados.
\end{procedimiento}

Las recomendaciones de este capitulo combinan tres niveles: comportamiento implementado por la plataforma, practica operativa recomendada y politica institucional pendiente. Cuando una decision depende de la organizacion, se marca como \campointerfaz{PENDIENTE DE DEFINICION INSTITUCIONAL}.

\section{Antes del torneo}

\begin{longtable}{p{3.4cm}p{7.6cm}p{2.0cm}}
\toprule
Revision & Recomendacion & Tipo\\
\midrule
\endhead
Acceso & Inicie sesion con el usuario que operara el panel y verifique que puede abrir resumen, divisiones y comunicaciones. & Implementado\\
Participantes & Revise registros, duplicados, categoria, institucion, robot e integrantes antes de sincronizar equipos. & Operativa\\
Asistencia & Envie solicitudes solo desde un entorno autorizado y revise confirmados antes de sincronizar. & Operativa\\
Equipos & Confirme que los equipos aprobados aparecen en la division correcta. & Operativa\\
Divisiones & Verifique colegios y universidades por separado; no mezcle decisiones entre categorias. & Implementado\\
Competencia & Genere o asegure tablero solo con datos revisados. No regenere con progreso sin entender el impacto. & Operativa\\
Navegador & Use un navegador actualizado y mantenga una sesion limpia para el operador principal. & Operativa\\
Red & Pruebe conectividad desde el lugar de operacion antes de iniciar resultados. & Operativa\\
Dispositivo alterno & Defina un equipo alterno con acceso autorizado por si el equipo principal falla. & Politica\\
Permisos & Identifique quien puede operar torneo, comunicaciones y Admin. & Politica\\
Responsables & Defina responsable de mesa, operador de plataforma y persona autorizada para correcciones. & Politica\\
Estado previo & Revise que no haya cambios pendientes ni datos demo mezclados con operacion real. & Operativa\\
Datos de demostracion & Elimine o aisle datos de demostracion segun procedimiento aprobado antes de trabajar con datos reales. & Politica\\
\bottomrule
\end{longtable}

\section{Durante el torneo}

\begin{itemize}
    \item Use una sola fuente de verdad para resultados: acta o decision oficial antes de guardar en la plataforma.
    \item Guarde y verifique. Despues de cada resultado sensible, confirme que la pantalla conserva el dato al recargar.
    \item Evite que varios operadores trabajen sobre la misma division sin coordinacion.
    \item Si aparece una confirmacion de correccion, cancele salvo que la mesa haya autorizado el cambio.
    \item Despues de una correccion, revise fases posteriores como indica \cref{cap:correcciones-manuales}.
    \item No use Django Admin cuando exista una interfaz operativa para la accion.
    \item Registre incidentes: hora aproximada, fase, accion, mensaje y decision tomada.
    \item Proteja datos personales en pantallas, capturas y conversaciones de soporte.
    \item Use \campointerfaz{Simulacion segura} para revisar comunicaciones durante pruebas; no ejecute envio real sin autorizacion.
\end{itemize}

\begin{operacionsensible}
Durante el torneo, detenerse para verificar es preferible a encadenar correcciones. La aplicacion no conserva cambios no guardados como mecanismo de contingencia; si hay duda, recargue y confirme el estado persistido.
\end{operacionsensible}

\section{Despues del torneo}

\begin{tabularx}{\linewidth}{>{\bfseries}p{4cm}X}
\toprule
Actividad & Recomendacion\\
\midrule
Podio & Revise campeon, segundo y tercer lugar si aplica. Confirme que el podio fue guardado.\\
Resultados finales & Verifique grupos, batallas, repechaje y final antes de declarar cierre operativo.\\
Comunicaciones & Consulte historial de comunicaciones para simulado, prueba, enviado, omitido o fallido.\\
Sesiones & Cierre sesiones de panel interno y Admin en equipos compartidos.\\
Evidencia & Conserve solo capturas o reportes autorizados y sin datos sensibles innecesarios.\\
Respaldo & Solicite o ejecute respaldo solo segun politica aprobada por responsable autorizado.\\
Incidentes & Documente errores, correcciones, envios fallidos y decisiones manuales.\\
Cambios posteriores & Evite modificar resultados o registros despues del cierre sin autorizacion organizativa.\\
\bottomrule
\end{tabularx}

\section{Plan de contingencia}

\begin{longtable}{p{3.6cm}p{5.9cm}p{3.6cm}}
\toprule
Incidente & Accion inmediata & Que no hacer\\
\midrule
\endhead
Perdida temporal de red & Detenga la operacion, no cierre conclusiones y verifique estado al volver la conexion. & No asumir que el ultimo cambio quedo guardado.\\
Cierre accidental del navegador & Reabra la pantalla, inicie sesion si hace falta y revise el ultimo estado persistido. & No repetir resultados sin comparar con acta.\\
Recarga inesperada & Verifique si la seleccion sigue visible y si hubo mensaje de exito previo. & No avanzar de fase por memoria.\\
Error de servidor & Capture mensaje sin datos sensibles y reporte si se repite. & No usar Admin para forzar cambios sin diagnostico.\\
Duda sobre un resultado & Detenga guardado y consulte la fuente oficial. & No aceptar confirmaciones sensibles.\\
Conflicto entre operadores & Pause la division, defina un operador responsable y compare estado persistido. & No operar simultaneamente la misma fase.\\
Informacion visual desactualizada & Recargue y compare con historial o pantalla relacionada. & No crear duplicados.\\
No es posible continuar fase & Revise resultados pendientes, cierre de fase y mensajes de bloqueo. & No crear fase manual para ocultar resultados incompletos.\\
\bottomrule
\end{longtable}

\section{Politicas pendientes}

Los siguientes puntos dependen de la organizacion y no de una funcion automatica descrita en este manual:
\begin{itemize}
    \item Canal formal de soporte: \campointerfaz{PENDIENTE DE DEFINICION INSTITUCIONAL}.
    \item Responsable autorizado para correcciones sensibles.
    \item Politica de respaldo antes y despues del torneo.
    \item Politica de conservacion de evidencias y capturas.
    \item Procedimiento institucional para retirar datos demo antes de produccion.
    \item Criterios para envio real de comunicaciones.
\end{itemize}

\section{Resumen para mesa de operacion}

\begin{itemize}
    \item Preparar: acceso, participantes, asistencia, equipos, divisiones y conectividad.
    \item Operar: seleccionar, guardar, verificar y avanzar solo con resultados completos.
    \item Corregir: detenerse, autorizar, aplicar y revisar fases posteriores.
    \item Comunicar: simular primero, probar de forma controlada y enviar real solo con aprobacion.
    \item Cerrar: revisar podio, historial, incidentes, sesiones y pendientes.
\end{itemize}
